\documentclass[12pt,a4paper]{article}

% ============================================
% PACKAGES
% ============================================
\usepackage[utf8]{inputenc}
\usepackage[english]{babel}
\usepackage{graphicx}
\usepackage{amsmath}
\usepackage{amssymb}
\usepackage{hyperref}
\usepackage{cite}
\usepackage{url}
\usepackage{booktabs}
\usepackage{array}
\usepackage{multirow}
\usepackage{subcaption}
\usepackage{xcolor}
\usepackage{listings}
\usepackage[margin=1in]{geometry}

% ============================================
% HYPERREF SETTINGS
% ============================================
\hypersetup{
    colorlinks=true,
    linkcolor=blue,
    filecolor=magenta,      
    urlcolor=cyan,
    citecolor=blue,
}

% ============================================
% CODE LISTING SETTINGS
% ============================================
\lstset{
    basicstyle=\ttfamily\small,
    breaklines=true,
    frame=single,
    language=Python
}

% ============================================
% TITLE INFORMATION
% ============================================
\title{\textbf{Synthetic-to-Real Image Translation\\for Chessboard Rendering}}

\author{
    Your Name \and 
    Partner Name \\
    \\
    Department of Computer Science, [University Name]\\
    \texttt{\{email1, email2\}@university.edu}
}

\date{\today}

% ============================================
% DOCUMENT
% ============================================
\begin{document}

\maketitle

% ============================================
% ABSTRACT
% ============================================
\begin{abstract}
In this project, we address the challenge of Synthetic-to-Real image translation for chessboard rendering. Our goal is to generate realistic chessboard images from synthetic layouts while preserving geometric consistency. We propose a robust pipeline utilizing the \textbf{Brownian Bridge Diffusion Model (BBDM)} integrated with a fine-tuned \textbf{VQGAN} for latent space representation. To ensure high-quality training data, we developed a preprocessing pipeline using \textbf{ffmpeg} for frame deduplication and the \textbf{Segment Anything Model (SAM)} for removing occluded samples (e.g., hands). While initial experiments with GAN-based methods (e.g., CUT) yielded suboptimal results, our diffusion-based approach demonstrates superior visual fidelity. We present quantitative and qualitative results, including training progression across epochs, demonstrating the effectiveness of our method.
\end{abstract}

% ============================================
% 1. INTRODUCTION
% ============================================
\section{Introduction}
\label{sec:introduction}

\subsection{Task Description and Motivation}
The task involves translating synthetic chessboard images—generated via rendering engines—into photorealistic images that mimic real-world capture conditions. This is crucial for bridging the domain gap in computer vision tasks where real labeled data is scarce but synthetic data is abundant.

\subsection{Challenges and Goals}
Key challenges include:
\begin{itemize}
    \item \textbf{Geometric Preservation:} The translation must not alter the position or identity of chess pieces.
    \item \textbf{Data Quality:} Real-world data often contains artifacts (motion blur, occlusions by hands) that can degrade model performance.
    \item \textbf{Texture Realism:} Accurately synthesizing lighting, shadows, and material properties.
\end{itemize}

\subsection{Main Contributions}
Our main contributions are:
\begin{enumerate}
    \item A comprehensive data curation pipeline using \textbf{ffmpeg} for temporal filtering and \textbf{SAM} \cite{sam2023} for occlusion removal.
    \item An enhanced synthetic data generation process using \textbf{Blender} with improved diversity (angles, lighting).
    \item Implementation and fine-tuning of a \textbf{BBDM} \cite{bbdm2023} architecture coupled with \textbf{VQGAN} \cite{vqgan2021} for high-fidelity image-to-image translation.
    \item A customized loss function strategy tailored for chessboard structure preservation.
\end{enumerate}

% ============================================
% 2. RELATED WORK
% ============================================
\section{Related Work}
\label{sec:related}

\subsection{Generative Adversarial Networks (GANs)}
GANs have been the standard for Image-to-Image (I2I) translation. Notable models include \textbf{Pix2Pix} for paired data and \textbf{CycleGAN} \cite{cyclegan2017} or \textbf{CUT} (Contrastive Unpaired Translation) \cite{cut2020} for unpaired data. We initially explored these methods but found them limited in handling the complex texture mapping required for this task.

\subsection{Vector Quantized Models (VQGAN)}
\textbf{VQGAN} \cite{vqgan2021} combines convolutional approaches with transformers by learning a discrete codebook of image constituents. It serves as a powerful perceptual compression stage, allowing subsequent models to operate in a lower-dimensional latent space without losing high-frequency details.

\subsection{Diffusion Models and BBDM}
Diffusion models, such as DDPM \cite{ddpm2020}, have recently surpassed GANs in generation quality. Specifically, \textbf{BBDM} \cite{bbdm2023} models the translation between two domains as a stochastic process governed by a Brownian Bridge. This allows for direct image-to-image translation without requiring cycle consistency, making it highly suitable for paired synthetic-to-real tasks.

\subsection{Segmentation Foundation Models (SAM)}
The \textbf{Segment Anything Model (SAM)} \cite{sam2023} has revolutionized semantic segmentation. We leverage SAM not for generation, but as a robust filtering tool to detect and remove occlusions (such as players' hands) from our training dataset.

% ============================================
% 3. METHOD
% ============================================
\section{Method}
\label{sec:method}

\subsection{Data Generation and Preprocessing}
\label{subsec:data_gen}
A critical part of our approach was ensuring high-quality input data. We employed a three-stage pipeline:

\subsubsection{Temporal Filtering with ffmpeg}
The raw data consisted of video recordings (PGN-labeled games). Since chess positions remain static for multiple seconds, extracting every frame results in high redundancy. We used \textbf{ffmpeg} to detect identical consecutive frames and extract only unique stable frames, significantly reducing dataset size while maximizing diversity.

\subsubsection{Occlusion Removal with SAM}
Real-world chess videos often contain hand movements during piece manipulation. Training on occluded images introduces artifacts. We utilized \textbf{SAM} \cite{sam2023} to detect "person" or "hand" masks. Images containing significant occlusions were automatically discarded from the training set to ensure the model learns only clean board representations.

\subsubsection{Enhanced Synthetic Data (Blender)}
We utilized the provided Blender script but introduced several improvements to bridge the domain gap:
\begin{itemize}
    \item Randomized camera angles and focal lengths to match real-world variance.
    \item Improved lighting conditions and shadows.
    \item Color grading adjustments to better align with the specific color temperature of the real dataset.
\end{itemize}

\subsection{Model Architecture: BBDM with VQGAN}
We selected the \textbf{Brownian Bridge Diffusion Model (BBDM)} \cite{bbdm2023} as our core architecture.
\begin{itemize}
    \item \textbf{Latent Space (VQGAN):} Instead of operating in pixel space, we use a VQGAN \cite{vqgan2021} to compress images into a discrete latent representation. This reduces computational complexity and focuses the diffusion process on semantic structure.
    \item \textbf{Fine-Tuning Strategy:} We performed fine-tuning on two levels:
    \begin{enumerate}
        \item \textbf{VQGAN Fine-tuning:} Adapted the codebook to specifically capture chessboard textures and piece geometries.
        \item \textbf{BBDM Fine-tuning:} Trained the diffusion bridge specifically on the mapping between our synthetic renders and the cleaned real dataset.
    \end{enumerate}
\end{itemize}

\subsection{Loss Functions}
\label{subsec:loss}
To ensure both perceptual quality and structural fidelity, we optimized a composite loss function.
\begin{equation}
    \mathcal{L}_{total} = \lambda_{diff} \mathcal{L}_{diffusion} + \lambda_{perc} \mathcal{L}_{LPIPS} + \lambda_{struct} \mathcal{L}_{structure}
\end{equation}
We incorporated the LPIPS metric \cite{lpips2018} to improve perceptual similarity.

% ============================================
% 4. EXPERIMENTS
% ============================================
\section{Experiments}
\label{sec:experiments}

\subsection{Experimental Setup}
\begin{itemize}
    \item \textbf{Dataset:} [Number] pairs of Synthetic-Real images after filtering.
    \item \textbf{Hardware:} Trained on [GPU Name, e.g., NVIDIA A100/T4].
    \item \textbf{Hyperparameters:} Batch size: [X], Learning rate: [Y], Epochs: [Z].
\end{itemize}

\subsection{Qualitative Results}
We present the progression of our model throughout the training process. Figure \ref{fig:epochs} demonstrates how the model gradually learns to synthesize realistic textures.

\begin{figure}[ht]
    \centering
    \begin{tabular}{ccc}
        \includegraphics[width=0.3\linewidth]{example-image-a} & \includegraphics[width=0.3\linewidth]{example-image-b} & \includegraphics[width=0.3\linewidth]{example-image-c} \\
        (a) Synthetic Input & (b) Epoch 10 & (c) Epoch 50 \\
        \includegraphics[width=0.3\linewidth]{example-image-a} & \includegraphics[width=0.3\linewidth]{example-image-b} & \includegraphics[width=0.3\linewidth]{example-image-c} \\
        (d) Epoch 100 & (e) Final Output & (f) Real Reference
    \end{tabular}
    \caption{Training progression: Visualizing the domain translation at different epochs.}
    \label{fig:epochs}
\end{figure}

\subsection{Quantitative Evaluation}
\textit{TODO: Add your FID/LPIPS scores here.}

% ============================================
% 5. WHAT DID NOT WORK
% ============================================
\section{What Did Not Work}
\label{sec:not_work}

Prior to adopting BBDM, we experimented extensively with GAN-based approaches, specifically \textbf{CUT} \cite{cut2020}.

\begin{itemize}
    \item \textbf{Issue:} The GAN models struggled to maintain the precise geometry of chess pieces, often hallucinating pieces or distorting the board grid.
    \item \textbf{Result:} As seen in Figure \ref{fig:gan_fail}, the outputs suffered from mode collapse and unrealistic artifacts.
    \item \textbf{Conclusion:} This led us to switch to Diffusion Models, which offer more stable training and better structural preservation.
\end{itemize}

\begin{figure}[ht]
    \centering
    % \includegraphics[width=0.5\linewidth]{path_to_failed_gan_image.png}
    \fbox{\parbox{6cm}{\centering \vspace{1cm} Placeholder for Failed GAN Result \vspace{1cm}}}
    \caption{Example of failed translation using CUT/GAN.}
    \label{fig:gan_fail}
\end{figure}

% ============================================
% 6. ABLATION STUDY
% ============================================
\section{Ablation Study}
\label{sec:ablation}

To validate our pipeline components, we analyzed the impact of:
\begin{enumerate}
    \item \textbf{w/o SAM Filtering:} Training on raw data (including hands) resulted in the model generating ghostly artifacts over the board.
    \item \textbf{w/o VQGAN Fine-tuning:} Using a generic pre-trained VQGAN resulted in blurry piece edges.
\end{enumerate}

\begin{table}[ht]
\centering
\caption{Ablation Study Results}
\begin{tabular}{lc}
\toprule
\textbf{Configuration} & \textbf{Visual Observation} \\
\midrule
Full Pipeline (Ours) & Sharp, Realistic, Clean \\
w/o SAM Filtering & Artifacts, Occlusions \\
w/o VQGAN FT & Blurry Textures \\
\bottomrule
\end{tabular}
\end{table}

% ============================================
% 7. DISCUSSION & LIMITATIONS
% ============================================
\section{Discussion and Limitations}
While our BBDM approach yields high-quality results, the inference time is slower compared to single-pass GANs due to the iterative diffusion process. Future work could explore distillation techniques to speed up sampling.

% ============================================
% 8. CONCLUSION
% ============================================
\section{Conclusion}
We presented a robust Synthetic-to-Real translation pipeline for chessboard images. By combining careful data curation (SAM, ffmpeg) with advanced diffusion models (BBDM+VQGAN), we achieved significant improvements over traditional GAN baselines.

% ============================================
% REFERENCES
% ============================================
\begin{thebibliography}{99}

% 1. CUT
\bibitem{cut2020}
Park, T., Efros, A. A., Zhang, R., \& Zhu, J. Y. (2020).
\newblock Contrastive learning for unpaired image-to-image translation.
\newblock In \textit{European Conference on Computer Vision} (pp. 319-345). Springer, Cham.

% 2. BBDM
\bibitem{bbdm2023}
Li, B., Xue, K., Liu, B., \& Lai, Y. K. (2023).
\newblock BBDM: Image-to-image Translation with Brownian Bridge Diffusion Models.
\newblock In \textit{Proceedings of the IEEE/CVF Conference on Computer Vision and Pattern Recognition} (pp. 1952-1961).

% 3. CycleGAN
\bibitem{cyclegan2017}
Zhu, J. Y., Park, T., Isola, P., \& Efros, A. A. (2017).
\newblock Unpaired image-to-image translation using cycle-consistent adversarial networks.
\newblock In \textit{Proceedings of the IEEE international conference on computer vision} (pp. 2223-2232).

% 4. VQGAN
\bibitem{vqgan2021}
Esser, P., Rombach, R., \& Ommer, B. (2021).
\newblock Taming transformers for high-resolution image synthesis.
\newblock In \textit{Proceedings of the IEEE/CVF conference on computer vision and pattern recognition} (pp. 12873-12883).

% 5. SAM
\bibitem{sam2023}
Kirillov, A., Mintun, E., Ravi, N., Mao, H., Rolland, C., Gustafson, L., ... \& Girshick, R. (2023).
\newblock Segment anything.
\newblock \textit{arXiv preprint arXiv:2304.02643}.

% --- SUGGESTED REFERENCES ---

% LPIPS (Highly recommended as you mentioned loss functions)
\bibitem{lpips2018}
Zhang, R., Isola, P., Efros, A. A., Shechtman, E., \& Wang, O. (2018).
\newblock The unreasonable effectiveness of deep features as a perceptual metric.
\newblock In \textit{Proceedings of the IEEE conference on computer vision and pattern recognition} (pp. 586-595).

% DDPM (Standard citation for diffusion background)
\bibitem{ddpm2020}
Ho, J., Jain, A., \& Abbeel, P. (2020).
\newblock Denoising diffusion probabilistic models.
\newblock \textit{Advances in Neural Information Processing Systems}, 33, 6840-6851.

\end{thebibliography}

\end{document}